\documentclass{article}
\usepackage{amsmath, amssymb, amsthm}
\usepackage{hyperref}
\usepackage{geometry}
\geometry{margin=1in}

\title{Erd\H{o}s Problem \#326}
\date{Last updated: 2025-08-31}
\author{Source: erdosproblems.com}

\begin{document}
\maketitle

\noindent\textbf{Status:} OPEN \quad \textbf{No prize}

\noindent\textbf{Tags:} number theory, additive basis

\medskip

\noindent\textbf{Problem Statement:}

Does there exist $A=\{a_1&#60;a_2&#60;\cdots\}\subset \mathbb{N}$ which is a minimal basis of order $2$ (i.e. every large integer is the sum of $2$ elements from $A$, and no proper subset of $A$ has this property), such that\[\lim_{k\to \infty}\frac{a_k}{k^2}=c\]for some $c\neq 0$?




Erd\H{o}s originally asked this for any basis (not necessarily minimal); such a basis was constructed by Cassels \cite{Ca57}. Erd\H{o}s and Graham conjectured a negative answer to this question.




References

[Ca57] Cassels, J. W. S., \"{U}ber Basen der nat\"{u}rlichen Zahlenreihe . Abh. Math. Sem. Univ. Hamburg (1957), 247-257.


\bigskip
\noindent\small{Source: \url{https://www.erdosproblems.com/326}}
\end{document}
